\cotefile{29}
\scanpage{22}
% archive sheet number: 21
sheaf associated to the preceding presheaf (NB the hoax-like condition imposed on the $(D_i)$ ensures that we do indeed have a presheaf). Then locally for the étale topology, we therefore have a situation $(D_i)$ as in a). This suffices to define a ramification datum $\underline{R}$ on $B$, moreover generalizing the one envisaged in a) (at least when the family $(D_i)_{i\quad I}$ of a) satisfies the additional local condition that we\handwritten{\unsure{g|?}} have just formulated). I see two special cases having some interest;. \deleted{\unsure{S|R|fi|if}} and \deleted{1\textsuperscript{\uline{o}})} where $\quad$ can be defined in terms of a divisor $D$ given on $S$:

1\textsuperscript{\uline{o}}) $S$ is locally noetherian, $D$ has normal crossings, we take as $D_i$ the regular local irreducible components of $D$ (which are defined only locally for the étale topology).

2\textsuperscript{\uline{o}}) $S$ is over a prescheme $T$, $D$ has normal crossings relative to $T$, we take as $D_i$ the local components smooth over $\unsure{\deleted{T}S|\deleted{S}T}$ (which are defined only locally for the étale topology).

In case 1\textsuperscript{\uline{o}}), the category $\operatorname{Rev}(_{\underline{R}}(S)$ is equivalent to a full subcategory of the category of coverings of $S$, and unless I am mistaken (it would need to be checked) the same holds in case 2\textsuperscript{\uline{o}}). The coverings in question are called, in the\deleted{ir} case 1\textsuperscript{\uline{o}}), tamely ramified coverings in the broad sense relative to the divisor $D$ (NB one indeed checks that this is a broader notion than that of ``tamely ramified covering with respect to $D$'', (meaning : with respect to the family reduced to $D$). Abhyankar's theorem allows these coverings to be characterized by the condition of being étale over \deleted{$S\quad\unsure{\underline{r}|\underline{a}|\underline{R}}$} $S-\operatorname{supp}D$, normal at the points over $\operatorname{supp}D$, every irreducible component of theirs d\unsure{\deleted{i}o|\deleted{o}i}minates one of $S$, finally they are tamely ramified in the classical sense at the maximal point\inserted{s} of $D$. In case 2\textsuperscript{\uline{o}}, it would be appropriate to call the coverings obtained ``tamely ramified for $D$, relative to the
